\section{Finite monoids, syntactic contexts, and linear grammars}
\label{sec:monoids}

A finite monoid is a supplied multiplication table $\mu:[m]^2\to[m]$,
$m\ge1$, and an element $e<m$, satisfying
\[
 \mu(e,x)=\mu(x,e)=x,\qquad
 \mu(\mu(x,y),z)=\mu(x,\mu(y,z)).
\]
These laws and table well-formedness are decidable bounded assertions. A
recognizer additionally supplies letter images $\lambda(a)<m$ and accepting
bits $P(x)$. Its word evaluation starts at $e$ and multiplies the letter
images from left to right. All structures here are finite numerical tables;
no quotient of an arbitrary set of words is assumed.

\begin{theorem}[Finite-monoid recognition and automata]\label{mon:recognition}
A language has a supplied finite-monoid recognizer if and only if it has a
finite DFA. Both conversions are primitive recursive. A monoid with $m$
elements gives an $m$-state DFA; an $n$-state DFA gives a monoid with at most
$n^n$ elements.
\end{theorem}
\begin{proof}
From a monoid take state set $[m]$, initial state $e$, transition
$x\mathrel{\smash{\xrightarrow{a}}}\mu(x,\lambda(a))$, and final bits $P$.
Its run is exactly the monoid evaluation. Associativity and the unit laws
also prove $\widehat\lambda(uv)=\mu(\widehat\lambda(u),\widehat\lambda(v))$
by induction along $v$.

Conversely represent a transformation of $[n]$ by its $n$-entry vector.
There are $n^n$ vectors. For transformations $f,g$, define
\[
 (f\cdot g)(q)=g(f(q));
\]
thus multiplication means first $f$, then $g$. The identity vector is a
two-sided unit, and both ways of multiplying three vectors have entry
$h(g(f(q)))$. This proves the monoid laws by a bounded comparison of entries.
Send $a$ to the transition vector $q\mapsto\delta(q,a)$. Run composition
shows that the image of $w$ is
$T_w(q)=\delta^*(q,w)$ and $T_{uv}=T_u\cdot T_v=T_v\circ T_u$.
Mark vector $f$ accepting exactly when $f(q_0)$ is an accepting state of
the DFA. This gives precisely the original accepting bit. Vector encoding,
composition, and the finite multiplication table are primitive recursive.
Their proofs compare computed entries, so the required inductions are
quantifier-free after bounded Boolean compilation.
\end{proof}

\begin{theorem}[Computing the generated transition monoid]\label{mon:generated}
From an $n$-state DFA one can compute its generated transition monoid,
its multiplication table, and a representative word for each element.
It has at most $N=n^n$ elements, and every representative may be chosen
with length below $N$.
\end{theorem}
\begin{proof}
Begin with the identity vector and repeatedly add $f\cdot T_a$ for every
currently present vector $f$ and letter $a$. After $j$ rounds the marked
vectors are exactly those represented by words of length at most $j$:
append one letter in the forward direction and remove the last letter in
the reverse direction. Store a first predecessor and letter for each new
vector. The universe has $N$ elements, so after $N$ rounds the bit table
is stable by bounded saturation. A shortest path from the identity to any
marked vector in this letter-transition graph has fewer than $N$ edges,
by finite path shortening. Those paths give the claimed representatives.

The stable set contains every word transformation: induction along a
supplied word uses its closure under each letter. It is closed under
multiplication as well. If $g=T_v$ is represented by its computed word,
starting at $f$ and reading $v$ remains in the stable set and ends at
$f\cdot g$. Thus the inherited table is a monoid. Renumber the marked
vectors in numerical order to obtain its finite code. Its recognizing
bits are inherited from Theorem~\ref{mon:recognition}. Completeness is
proved by the finite-table invariant and arbitrary-word induction, not
by extrapolating from sampled words.
\end{proof}

\begin{theorem}[Effective syntactic monoid and short distinguishing contexts]
\label{mon:syntactic}
Let $M$ be a reachable minimal $n$-state DFA for $L$. Define two-sided
syntactic equivalence by
\[
 u\approx_L v\quad\longleftrightarrow\quad
 \forall x,y\;\bigl(L(xuy)\leftrightarrow L(xvy)\bigr).
\]
Then $u\approx_L v$ if and only if $T_u=T_v$ as $n$-entry vectors.
When the vectors differ, one can compute distinguishing contexts with
$|x|<n$ and $|y|\le n^2$. Consequently the syntactic monoid of a supplied
regular language is effectively finite and primitive-recursively computable.
\end{theorem}
\begin{proof}
If the vectors agree, a prefix $x$ reaches some state $q$; reading $u$ or
$v$ then gives the same state, and the common suffix $y$ gives the same
accepting bit. This proves the context equation for arbitrary supplied
$x,y$.

If the vectors differ, bounded search finds $q<n$ with
$p=T_u(q)\ne T_v(q)=r$. Reachability provides an access word $x$ for $q$
with $|x|<n$. Distinct states of a minimal DFA are distinguishable;
otherwise the quotient construction in Theorem~\ref{reg:minimize} would
reduce its number of states. The table of
Proposition~\ref{reg:distinguish} therefore provides $y$, $|y|\le n^2$,
whose accepting bits from $p,r$ differ. Hence $xuy,xvy$ distinguish the
contexts. Since vector equality is decidable, this proves both directions
constructively, including the implication from the universally quantified
context equation to vector equality.

Equivalence is a two-sided congruence because multiplication of equal
vectors on either side preserves equality. More explicitly,
$T_{aub}=T_a\cdot T_u\cdot T_b$ for arbitrary words $a,b$.
The generated transition monoid therefore represents exactly the syntactic
classes, with class multiplication represented by concatenation. Given any
input DFA, first apply the primitive recursive minimization procedure,
then Theorem~\ref{mon:generated}. This returns the actual finite multiplication
table and representatives, without forming an infinite quotient. The
unbounded context quantifiers occur only in the theorem statement;
its test is equality of finitely many entries, and its negative branch
supplies the explicit bounded contexts above.
\end{proof}

\begin{proposition}[Left-linear grammars and the danger of mixing directions]
\label{mon:leftlinear}
Finite left-linear grammars, with productions $A\to Bu$ or $A\to u$,
define exactly the regular languages. Allowing both left-linear and
right-linear productions in one grammar can define a nonregular language.
\end{proposition}
\begin{proof}
Replace each left-linear production $A\to Bu$ by $A\to u^RB$, and each
$A\to u$ by $A\to u^R$. A left-linear sentential form consists of one
nonterminal followed by a terminal suffix. Reversing it gives a terminal
prefix followed by that nonterminal, and each old production corresponds
to exactly one new production. Induction along the supplied derivation
therefore identifies the new grammar's language with the reversal of the
old one's language. The new grammar is right-linear, so
Theorem~\ref{reg:grammar} and closure under reversal give a DFA for the
old grammar. Conversely convert a DFA for $L^R$ to a right-linear grammar
and reverse this construction to obtain a left-linear grammar for $L$.
All reversals and derivation transformations are finite PR scans.

For a strict separation, use
\[
 S\to aA\mid\eps,\qquad A\to Sb.
\]
Every nonterminal derivation alternates the sentential forms
$a^jSb^j$ and $a^{j+1}Ab^j$. Choosing $S\to\eps$ terminates with
$a^jb^j$; choosing the other rule and then the forced $A$ rule increments
$j$. This bounded trace invariant proves that the language is exactly
$\{a^jb^j:j\in\N\}$, with every such word obtained by $j$ iterations
and the final empty production. It is nonregular by
Corollary~\ref{reg:nonregular}. Each production is individually linear,
but their two directions cannot be mixed while retaining the regular-language
characterization.
\end{proof}
